\bibitem{hayesroth1985} B. Hayes-Roth. A blackboard architecture for control. \emph{Artificial Intelligence} 26:251--321, 1985. doi:10.1016/0004-3702(85)90063-3.

\bibitem{nii1986} H. P. Nii. The Blackboard Model of Problem Solving and the Evolution of Blackboard Architectures. \emph{AI Magazine} 7(2):38--53, 1986.

\bibitem{baars1988} B. J. Baars. \emph{A Cognitive Theory of Consciousness}. Cambridge University Press, 1988.

\bibitem{dehaene2001} S. Dehaene and L. Naccache. Towards a cognitive neuroscience of consciousness: basic evidence and a workspace framework. \emph{Cognition} 79:1--37, 2001. PMID:11164022.

\bibitem{mashour2020} G. A. Mashour, P. Roelfsema, J.-P. Changeux, and S. Dehaene. Conscious Processing and the Global Neuronal Workspace Hypothesis. \emph{Neuron} 105:776--798, 2020. doi:10.1016/j.neuron.2020.01.026.

\bibitem{bengio2017} Y. Bengio. The Consciousness Prior. arXiv:1709.08568, 2017.

\bibitem{goyal2021} A. Goyal et al. Coordination Among Neural Modules Through a Shared Global Workspace. arXiv:2103.01197, 2021.

\bibitem{gurnee2026} W. Gurnee et al. Verbalizable Representations Form a Global Workspace in Language Models. \emph{Transformer Circuits Thread}, published 6 July 2026; arXiv:2607.15495, posted 16 July 2026.

\bibitem{furnas1987} G. W. Furnas, T. K. Landauer, L. M. Gomez, and S. T. Dumais. The Vocabulary Problem in Human-System Communication. \emph{Communications of the ACM} 30(11):964--971, 1987. doi:10.1145/32206.32212.

\bibitem{swanson1986} D. R. Swanson. Fish oil, Raynaud's syndrome, and undiscovered public knowledge. \emph{Perspectives in Biology and Medicine} 30:7--18, 1986. doi:10.1353/pbm.1986.0087.

\bibitem{garikaparthi2025} A. Garikaparthi, M. Patwardhan, A. S. Kanade, A. Hassan, L. Vig, and A. Cohan. MIR: Methodology Inspiration Retrieval for Scientific Research Problems. In \emph{Proceedings of ACL 2025}, pages 28614--28659. doi:10.18653/v1/2025.acl-long.1390.

\bibitem{carbonell1998} J. Carbonell and J. Goldstein. The Use of MMR, Diversity-Based Reranking for Reordering Documents and Producing Summaries. In \emph{Proceedings of SIGIR 1998}, pages 335--336, 1998. doi:10.1145/290941.291025.

\bibitem{jafarzadeh2021} M. Jafarzadeh et al. A Review of Open World Learning and Steps Toward Open World Learning Without Labels. arXiv:2011.12906, 2021.

\bibitem{saunders2018} B. Saunders et al. Saturation in qualitative research: exploring its conceptualization and operationalization. \emph{Quality \& Quantity} 52:1893--1907, 2018. doi:10.1007/s11135-017-0574-8.

\bibitem{tate2009} R. Tate, M. Stepp, Z. Tatlock, and S. Lerner. Equality Saturation: A New Approach to Optimization. In \emph{Proceedings of POPL 2009}, pages 264--276, 2009. doi:10.1145/1480881.1480915.

\bibitem{willsey2021} M. Willsey et al. egg: Fast and Extensible Equality Saturation. \emph{Proceedings of the ACM on Programming Languages} 5(POPL), 2021. doi:10.1145/3434304.

\bibitem{vanemden1976} M. H. van Emden and R. A. Kowalski. The Semantics of Predicate Logic as a Programming Language. \emph{Journal of the ACM} 23(4):733--742, 1976. doi:10.1145/321978.321991.

\bibitem{rissanen1978} J. Rissanen. Modeling by shortest data description. \emph{Automatica} 14(5):465--471, 1978. doi:10.1016/0005-1098(78)90005-5.

\bibitem{tishby2000} N. Tishby, F. C. Pereira, and W. Bialek. The information bottleneck method. arXiv:physics/0004057, 2000.

\bibitem{fletcher2026} A. H. A. Fletcher and M. Stevenson. Decision-Theoretic Stopping Rules for Document Screening. arXiv:2606.07071, 2026.

\bibitem{fletcher2026b} A. H. A. Fletcher and M. Stevenson. Confidence-Based Stopping Methods for Systematic Reviews. arXiv:2606.15380, 2026.

\bibitem{macfarlane2022} A. MacFarlane, T. Russell-Rose, and F. Shokraneh. Search strategy formulation for systematic reviews: Issues, challenges and opportunities. \emph{Intelligent Systems with Applications}, 2022. doi:10.1016/j.iswa.2022.200091.

\bibitem{shi2017} B. Shi and T. Weninger. Open-World Knowledge Graph Completion. arXiv:1711.03438, 2017.

\bibitem{yang2022} H. Yang, Z. Lin, and M. Zhang. Rethinking Knowledge Graph Evaluation Under the Open-World Assumption. arXiv:2209.08858, 2022.

\bibitem{li2024kg} J. Li, R. Luo, J. Sun, J. Xiao, and Y. Yang. Progressive Knowledge Graph Completion. arXiv:2404.09897, 2024.

\bibitem{gradel2024} E. Gr\"adel and V. Tannen. Provenance Analysis and Semiring Semantics for First-Order Logic. arXiv:2412.07986, 2024.
